\scanpage{20}

\uline{Géom. Algébrique}

Compléments sur la formation des $\varprojlim$ et des sommes dans $\mathrm{Aff}_k$.

a) \uline{Les sommes \inserted{finies} sont disjointes et universelles}

\[
X=\coprod_{i\in I}X_i\longleftarrow X_i
\]

1) $X_i\to X$ un mono et $=$ une \uline{imm. fermée} \qquad ($A=\prod A_i\to A_i$ surjectif)

2) $X_i\times_X X_j=\varnothing$ si $i\ne j$ \qquad i.e. $A_i\otimes_A A_j=0$ si $i\ne j$ car l'\unsure{idemt|élément|égal} $1$ de $A_i$, $0$ de $A_j$ \deleted{\unsure{égal à $1$ dans $A_j$, $0$ dans $A_i$}}

\diagram{\begin{tikzcd}[ampersand replacement=\&,column sep=normal,row sep=small]
3)\quad\forall\quad X'\arrow[r]\&X\&\text{on a}\quad X'\xleftarrow{\sim}\coprod X'_i\\
X'_i\arrow[u]\arrow[r]\&X_i\arrow[u]\&\unsure{q.}
\end{tikzcd}}

i.e. $\forall\,A=\prod\unsure{A_i|A_j}\to B$, $B$ décomposé en $\prod B_i$, et \ldots

\deleted{\unsure{Donc dans ce cas|b) Dans le cas où}} Donc les produits \inserted{(ou \uline{produits fibrés})} sont distributifs p. r. aux sommes

\[
\left(\coprod X_i\right)\times\left(\coprod_j Y_j\right)\simeq\coprod_{ij}(X_i\times Y_j)
\]

\[
\bigl[\,\deleted{\unsure{Y}}X\times\underbrace{\coprod Y_i}_{Y}\longleftarrow\coprod_i X\times Y_i\qquad\ldots\,\bigr]
\]

\diagram{\begin{tikzcd}[ampersand replacement=\&,row sep=small]
X\times Y_i\arrow[r]\arrow[d]\&X\times Y\arrow[d]\&X\arrow[d,dash]\\
Y_i\arrow[r]\&Y\arrow[r,dash]\&e
\end{tikzcd}}

Donc le foncteur

\[
\begin{aligned}
(\mathrm{Ens}\,\unsure{\mathrm{fin}|\mathrm f|\mathrm{finis}})&\longrightarrow\mathrm{Aff}_k\\
\underline I&\longrightarrow I_k=\coprod_I e_k
\end{aligned}
\]

commute aux \deleted{$\unsure{\text{limites}}\ \varprojlim$ finies} produits \inserted{finis}, \deleted{\unsure{en particulier aux}} \inserted{et $=$ aux} produits fibrés, donc aux $\varprojlim$ finies. Donc transforme monoïde (groupe) fini en monoïde (groupe) algébrique affine.

Du point de vue alg. comm., on a

\margin{Le foncteur $I_k$ est le foncteur : ens des décompositions idempotentes indexées par $I$}

\[
I_k=V_{k^I}\qquad k\to k^I\quad\text{contravariant en $I$}
\]

\inserted{algèbre produit $\simeq\mathrm{Fonc}(I,k)$}

$k^I\simeq k[I]$ \inserted{(pr. finis)} (gare aux confusions avec polynômes, Si \deleted{\unsure{$I$ a}} \inserted{$I$ a} une \deleted{\unsure{grpe, la multiplication}} \inserted{loi de composition} interne $I\times I\to I$, cela définit

\[
k[I]\longrightarrow k[I\times I]\simeq k[I]\otimes k[I]
\]

la comultiplication \ldots.

% archive sheet number: 19
